% Galen Phase 2 preprint — ACL-style LaTeX (arXiv-ready)
% Engine: XeLaTeX (Unicode + polytonic Greek via Times New Roman)
% Build: see README.md
\documentclass[11pt]{article}

% preprint = non-anonymous, with page numbers (arXiv-friendly)
\usepackage[preprint]{acl}

\usepackage{fontspec}
\usepackage{microtype}
\usepackage{graphicx}
\usepackage{booktabs}
\usepackage{xcolor}
\usepackage{enumitem}

% Main font: Windows Times New Roman covers polytonic Greek (Unicode).
% Greek is entered as Unicode; no separate babel language switch required.
\setmainfont{Times New Roman}
\setmonofont{Consolas}[Scale=0.9]

\title{Which Words Block the Machine Translation of Galen?\\
A Reproducible Terminology Density Map}

\author{William Pepperman \\
  Independent researcher · Montblanc Technologies \\
  \texttt{montblanctech@outlook.com}}

\begin{document}
\maketitle

\begin{abstract}
The obstacle to machine-translating the large untranslated Galenic corpus has recently been
\textit{measured}: a 2026 MQM study reports that passage-level terminology rarity predicts MT quality on
Galen's previously untranslated pharmacological work at $r = -.97$ --- the failure mode is vocabulary, not syntax. What that measurement does not
supply is a queryable answer to \textit{which} Galenic vocabulary is the blocker. We build that resource.
Over the digitized Galenic corpus (First1KGreek), we run a reproducible, config-driven pipeline ---
extraction, odyCy lemmatization, log-likelihood keyness against a 19.68M-token non-Galen Greek
reference, and LSJ cross-reference --- and introduce a \textbf{four-way classification} that separates
genuinely rare vocabulary from domain-specific and polysemous vocabulary, a distinction
whole-corpus frequency conflates. The classification isolates, \textit{a priori} from frequency alone,
a set of common-but-technical terms (οὐσία, δύναμις, ποιότης) that rarity-based difficulty predictors
cannot see by construction --- a candidate failure mode orthogonal to the rarity the MT literature
has measured. We release a prioritized punch list of ${\sim}$1,400 candidate terms, the full pipeline (MIT),
and all data (CC-BY-4.0), each stage independently re-runnable over a different subcorpus. We claim
an instrument, not a translation, and not lexicon adjudication.
\end{abstract}

\section{Introduction}
\label{sec:intro}

A large fraction of the Galenic corpus --- the single biggest body of surviving ancient Greek prose ---
has never been translated into any modern language. The constraint is throughput, not willingness:
the specialist labor required to translate technical Greek medicine does not scale, and machine
translation has until recently been assumed unusable on it.

That assumption has now been quantified. A 2026 study applying MQM error annotation to
machine-translated Galen reports that passage-level \textbf{terminology rarity} predicts MT quality at
$r = -.97$ on the previously untranslated pharmacological work~\citep{galen-mt-mqm-2026}.
The practical consequence is a reframing: the task is not ``improve the model's grammar'' but ``identify the
vocabulary.'' If a small, well-chosen set of terms accounts for most of the degradation, then
supplying those terms as constrained vocabulary is a tractable intervention where full translation
is not.

No published resource says which terms those are. This paper builds the instrument panel: a
reproducible map of Galen's terminological density that tells an MT researcher, or a lexicographer,
\textit{where} the hard vocabulary concentrates and \textit{what kind} of hard each term is. We are explicit about
what this is not --- it is not a translation, and it is not an adjudication of what any term means.
Where a term's sense is contested, we surface the term and leave the field blank for a domain
expert to fill.

\paragraph{Contributions.}
\begin{enumerate}[leftmargin=*,itemsep=2pt]
\item A reproducible, config-driven pipeline over digitized Galen, from TEI extraction through a
  thin-gloss punch list, with the model and thresholds pinned.
\item A \textbf{four-way classification} of candidate-hard vocabulary along two frequency axes --- rarity in
  the medical subcorpus $\times$ rarity in the \textit{disjoint non-medical Galen} --- that distinguishes
  genuinely rare terms from domain-specific and polysemous ones.
\item A prioritized punch list of 1,433 candidate terms with attested example contexts, released as
  the scaffold for a future community-compiled Phase 3.
\item A demonstration that the frequency contrast partitions candidate-hard vocabulary, without
  semantic annotation, into two mechanistically distinct classes --- rarity-driven difficulty
  (which existing predictors capture) and common-but-sense-loaded vocabulary (which they cannot).
\end{enumerate}

\section{Background \& Related Work}
\label{sec:background}

\paragraph{Text source.}
We draw on the Open Greek and Latin / First1KGreek digitization of Galen~\citep{first1kgreek},
keyed by CTS URN. The Phase~1 census (\S\ref{sec:data}) fixes the work list and word counts and
is released separately~\citep{pepperman-2026-galen-census}
(DOI \url{https://doi.org/10.5281/zenodo.21561962}).

\paragraph{Greek NLP.}
We lemmatize with \textbf{odyCy}~\citep{kostkan-etal-2023-odycy}, the Aarhus Ancient-Greek
spaCy pipeline trained on Perseus and PROIEL treebanks, using the transformer model
\texttt{grc\_odycy\_joint\_trf}. odyCy is, to our knowledge, the strongest openly available lemmatizer for
ancient Greek; we cite it prominently because its accuracy is the pivot on which every downstream
frequency rests, and because its authors are the natural first readers of this work.

\paragraph{Keyness.}
We use log-likelihood (G$^2$) keyness in the standard corpus-linguistics form~\citep{dunning-1993-accurate,rayson-garside-2000-comparing},
comparing a target corpus against a reference corpus term by term.

\paragraph{Prior lexical work.}
The standard curated lexicon of Galenic medical vocabulary is Durling's \textit{A Dictionary of
Medical Terms in Galen}~\citep{durling-1993}. It is complementary to, not in competition with, this
work: Durling adjudicates what terms mean; we deliberately do not. Our map is \textit{a priori}
candidate generation from frequency structure --- it surfaces which vocabulary to examine and leaves
the sense to the lexicographer or the domain expert.

\paragraph{Relation to the motivating measurement.}
\citet{galen-mt-mqm-2026} is what makes this resource worth building, and it also fixes its scope.
Applying MQM error annotation to machine-translated Galen, they establish that passage-level
\textbf{terminology rarity} --- lemmas absent from, or infrequent in, a general literary Greek
reference --- predicts translation quality at $r = -.97$, with critical errors concentrated in rare
pharmacological terminology; conversely, they report that \textit{common} technical terms (e.g.\
δύναμις, κρᾶσις) are rendered without systematic difficulty. Their analysis is passage-level ---
twenty passages from two works --- and they explicitly invite extension to the wider corpus for the
statistical power a small sample cannot supply. We take up that extension and change the unit of
analysis from the passage to the term: rather than correlating a passage's rare-term ratio with its
score, we map, per lemma across the digitized corpus, \textit{where} the rare vocabulary concentrates
and --- via the medical-vs-non-medical contrast --- \textit{what kind} of difficulty each candidate
poses. Two consequences follow. First, the contrast isolates a class their rarity predictor cannot
reach by construction: common, sense-loaded vocabulary (\S\ref{sec:results}). Second, because both
their rare-term ratio and our counts rest on automatic lemmatization whose accuracy on Galenic
medical Greek has never been reported --- they lemmatize with Stanza, we with
odyCy~\citep{kostkan-etal-2023-odycy} --- we measure that accuracy directly (\S\ref{sec:accuracy})
rather than assume it. The \textbf{medical-vs-non-medical comparison within a single author} is, to
our knowledge, unpublished, and it is what makes the domain/polysemy distinction in
\S\ref{sec:results} possible.

\section{Data}
\label{sec:data}

\paragraph{Census.}
The Phase~1 census~\citep{pepperman-2026-galen-census}
(\url{https://doi.org/10.5281/zenodo.21561962}) is the ground-truth work list:
which editions exist, their URNs, category tags, and grc word counts. We treat it as authoritative
for \textit{which} works to process and do not re-derive it.

\paragraph{Subcorpora.}
We partition Galen into a \textbf{medical subcorpus} --- works tagged pharmacology,
therapeutics, pathology, pulse, semiotics, dietetics, anatomy, physiology --- and a \textbf{disjoint
non-medical subcorpus} (logic, protreptic, commentary, methodology, psychology, hygiene).
Pseudo-Galenic attributions are excluded via an explicit config flag, not a buried default, so the
choice stays auditable. Inclusion criteria are stated so a reviewer can contest any single
assignment.

\paragraph{Sizes.}

\begin{table}[t]
\centering
\begin{tabular}{lrr}
\toprule
\textbf{Corpus} & \textbf{Tokens} & \textbf{Works} \\
\midrule
Galen, medical & 1{,}746{,}283 & 71 \\
Galen, non-medical (disjoint) & 724{,}353 & 32 \\
Galen, whole & 2{,}470{,}636 & 103 \\
Non-Galen First1K reference & 19{,}680{,}000 ($\approx$) & --- \\
\bottomrule
\end{tabular}
\caption{Corpus sizes used in this study.}
\label{tab:sizes}
\end{table}

The same lemmatizer model is applied to every side. This is not a detail: keyness compares lemma
inventories across corpora, so a mixed-model comparison would score model differences as if they
were vocabulary differences.

\section{Method}
\label{sec:method}

The pipeline is a linear sequence of independently re-runnable stages; each reads one input and
writes one output, so any stage can be re-run without repeating the ones before it.

\begin{enumerate}[leftmargin=*,itemsep=3pt]
\item \textbf{Fetch} --- download the grc TEI editions named by the census; log 404s rather than failing.
\item \textbf{Extract} --- parse TEI, keep \texttt{xml:lang="grc"} edition text only, exclude apparatus (\texttt{<app>}/
  \texttt{<rdg>}) and restored/conjectural material. Normalize to Unicode NFC; map lunate sigma
  (ϲ/Ϲ $\rightarrow$ σ/ς, word-final ς); strip all-caps running headers; handle crasis/elision explicitly.
  A census-wordcount guardrail flags any work extracting below 70\% of its recorded count.
\item \textbf{Lemmatize} --- odyCy \texttt{grc\_odycy\_joint\_trf} (spaCy 3.7.5, spacy-transformers 1.3.5, model
  0.7.0, all pinned), chunked on sentence boundaries to avoid mid-sentence context loss.
  Per-token \texttt{(surface, lemma, pos, morph)} is retained.
\item \textbf{Frequency} --- aggregate to lemma $\rightarrow$ raw count and per-10{,}000-token normalized count, for each
  corpus separately.
\item \textbf{Keyness} --- log-likelihood G$^2$ for every medical-subcorpus lemma against the reference, with a
  minimum-domain-frequency floor so a lemma appearing twice cannot produce a spurious top score.
\item \textbf{LSJ cross-reference} --- look up each candidate in the machine-readable LSJ~\citep{lsj-perseus}. The
  Perseus LSJ index is keyed in Beta Code; we build a Unicode-NFC index once at load time (Beta
  Code $\rightarrow$ Unicode is the less ambiguous direction and yields a human-inspectable index) and record
  whether an entry exists and its \texttt{<sense>} count.
\item \textbf{Punch list} --- filter to thin-gloss candidates, attach example URNs and one attested context
  per term.
\end{enumerate}

\paragraph{The four-way classification.}
The core analytic step crosses two axes: rarity in the \textbf{medical
subcorpus} $\times$ rarity in the \textbf{disjoint non-medical Galen}. We use the disjoint non-medical corpus,
not whole-Galen, deliberately: the medical works are 71\% of the whole, so medical-frequency and
whole-frequency are collinear and the contrast collapses (a pure domain term like ἀμμωνιακός, absent
outside medicine, barely moves against the whole because the whole is mostly medicine). Against the
disjoint non-medical corpus the contrast is real: ἀμμωνιακός reads medical 1.40 / non-medical 0.00
per-10k. The rarity threshold is a config value; we report sensitivity to it (\S\ref{sec:results}). POS is restricted
to NOUN/ADJ/VERB.

\section{Lemmatizer Accuracy}
\label{sec:accuracy}

Every frequency in this paper inherits the lemmatizer's error rate, so we report it directly rather
than citing odyCy's published treebank scores --- Galen's pharmacological vocabulary is out-of-domain
relative to odyCy's training data, so those published numbers are an upper bound, not our number.

In the absence (to date) of a Greek-reading annotator, we bound the accuracy with three machine
measurements of different epistemic character --- a rule-based floor, a cross-system agreement
ceiling, and a gold-standard point estimate on the rare stratum --- rather than leaving it
unmeasured. All three are report-only: no output of the frozen pipeline is corrected in response.

\paragraph{Floor: morphological admissibility.}
For every token we ask whether odyCy's lemma appears among the readings that
Morpheus~\citep{crane-1991-morpheus} --- a rule-based analyzer sharing no training data with odyCy ---
admits for that surface form. On the medical subcorpus, \textbf{89.51\%} of analyzable tokens carry
an admissible lemma (4.41\% of tokens receive no Morpheus analysis and are excluded from the
denominator; they are reported, not silently pooled). Every inadmissible lemma is a near-certain
error, so this is a floor on correctness; it is a slightly permissive floor, because Morpheus itself
lists some variant stems as lemmas (it admits ἥπατος, a documented odyCy error for ἧπαρ, alongside
the correct reading).

\paragraph{Agreement: independent implementations, shared training data.}
Against Stanza's two Ancient Greek models~\citep{qi-etal-2020-stanza} on a 5{,}000-token random
sample, odyCy agrees with \texttt{grc\_perseus} pairwise on \textbf{82.89\%} of aligned tokens, and
with the two Stanza models on \textbf{94.17\%} of the tokens where those two agree with each other
--- a subset that conditioning selects for easiness. Both figures are optimistic bounds rather than
accuracy: all three systems train on overlapping Universal Dependencies treebanks, so their errors
are correlated, and 15.82\% of the sample was lost to tokenization misalignment. We report agreement
for what it is and use its disagreement set to target future human review where it is most
informative.

\paragraph{Rare-stratum gold standard: LSJ's own Galen citations.}
LSJ~\citep{lsj-perseus} cites Galen by K\"uhn volume and page, and each citation is a
lexicographer's judgment binding an attested form at a locus to a headword. We mapped these
citations onto the corpus through the editions' K\"uhn page breaks and accepted a token only when
exactly one surface form in the cited page window is morphologically derivable from the headword per
Morpheus (quoted dictionary text serves as context only, never as the scoring target, since LSJ
abbreviates headwords inside entries; the full location funnel, 7{,}706 citations $\rightarrow$
1{,}850 scored, ships with the data). On these \textbf{1{,}850} expert-attested loci, odyCy
reproduces the lexicographer's headword for \textbf{62.32\%} of tokens ignoring accentuation (Wilson
95\% CI 60.09--64.50), \textbf{53.95\%} under strict accent matching. LSJ-cited vocabulary is by
nature the rare, lexicographically interesting stratum, so this is a stratum accuracy, not an
overall one --- and read against the 89.51\% frequency-weighted floor, the two numbers measure, on
expert-annotated data and independently of any MT system, the same rarity-driven degradation that
\citet{galen-mt-mqm-2026} infer from translation error rates. The 697 mismatches decompose into:
identity passthrough of an inflected form (411), garbled non-word lemmas (135), substitution of a
wrong real lemma (136), and variant-headword conventions arguably not errors at all (15; scoring
these as correct moves the rate only to 63.14\%).

\paragraph{Machine-adjudicated estimate: the human protocol, piloted by two models.}
The salted-control protocol written for human readers (a frozen random 60-token sample, seed
20260725, shuffled with 13 known-answer control rows whose identity the reader is not told) was run
blind by two frontier language models in separate sessions
(\texttt{claude-fable-5-1} and \texttt{claude-opus-5-5}, 2026-11-05), each seeing only the
reviewer sheet. Both passed the control gate in full --- the planted non-word, all six documented
wrong lemmas and all five correct anchors caught, 13/13 each, including the two accent-only
controls --- which is the audit the protocol provides on the readers themselves. On the 60 random
tokens the readers marked \textbf{83.33\%} and \textbf{81.67\%} of lemmas correct (Wilson 95\% CI
71.97--90.69 and 70.08--89.44), with raw agreement 98.33\% and Cohen's $\kappa = 0.942$; of the ten
tokens both marked wrong, nine had already been flagged by the cross-system disagreement set above,
and nine of the ten corrections proposed by the two readers coincide (six are identity
passthrough of an inflected form, consistent with the rare-stratum taxonomy). We report this as a
\emph{machine-adjudicated estimate}, labelled as such and never pooled with the analyzer floor or the
lexicographer gold: a model's reading is a claim by a named model under a published prompt, not a
philologist's. It sits where the three machine bounds predicted --- between the frequency-weighted
floor and the rare-stratum gold --- and it makes the eventual human two-reader check, which remains
planned under the same protocol, cheaper and better targeted. The same reader then adjudicated,
system-blind, all 942 tokens on which the Stanza models disagreed with odyCy: it sided with odyCy
on 42.7\% overall (Wilson 95\% 39.6--45.9), but on the 479 disagreements that are not elision or
tokenization artifacts on only \textbf{14.0\%} (67/479; 11.2--17.4) --- so the cross-system
disagreement set is a strong error detector, and it is where a human reader should start. The
reviewer sheets, all answer files, the private keys and the scoring script ship with the data.

We treat the rare-stratum figure as the headline caveat on the punch list (\S\ref{sec:punch}),
whose candidates are drawn disproportionately from exactly that stratum.

\paragraph{Passthrough baseline (computed).}
As an independent check that Galen's vocabulary does not
degrade the lemmatizer relative to ordinary Greek, we measure content-word passthrough (the rate at
which a NOUN/VERB/ADJ/ADV token is emitted with \texttt{lemma == surface}, a proxy for the lemmatizer
falling back to the identity). On Galen this is \textbf{26.85\%}; on the 11.86M-token non-Galen reference
it is \textbf{26.41\%} ($\Delta$ $-0.44$ points). Passthrough is therefore odyCy's baseline behavior on Greek, not
Galen-specific degradation --- which pre-empts the natural objection that out-of-domain medical
vocabulary must be breaking the lemmatizer. It is not.

Residual, systematic error classes are carried transparently into the limitations (\S\ref{sec:limitations}) rather than
silently corrected, because correcting them is a lexical judgment we have chosen not to make.

\section{Results: The Four-Way Classification}
\label{sec:results}

Of 12{,}538 classified lemmas, at the threshold of 1.0 per-10k:

\begin{table}[t]
\centering
\begin{tabular}{lrl}
\toprule
\textbf{Quadrant} & \textbf{Count} & \textbf{Reading} \\
\midrule
\texttt{genuinely\_hard} (rare in both) & 11{,}709 & rare everywhere --- the candidate universe \\
\texttt{polysemy\_risk} (common in both) & 496 & common words carrying loaded technical senses \\
\texttt{false\_alarm} / domain (medical-only) & 190 & rare corpus-wide, at home in medicine \\
\texttt{mostly\_irrelevant} (non-medical-only) & 143 & not the medical vocabulary we are after \\
\bottomrule
\end{tabular}
\caption{Four-way classification at threshold 1.0 per-10k ($N=12{,}538$ lemmas).}
\label{tab:fourway}
\end{table}

The off-diagonal quadrants are where the analysis earns its keep, and they survive threshold
perturbation (at 0.5 / 1.5 per-10k the polysemy quadrant holds 853 / 348 members; the domain
quadrant 359 / 118 --- the structure is stable, not a threshold artifact).

\paragraph{The polysemy quadrant is what the contrast is for.}
Its members --- δύναμις, οὐσία, ποιότης, φύσις, σῶμα, λόγος, πάθος, εἶδος, ὕλη --- are high-frequency
on both sides: common words that also carry Galen-technical senses. They are, by construction,
invisible to any rarity signal. A term this frequent scores as \textit{easy} under corpus-frequency
rarity --- which is exactly the predictor the MT study operationalizes~\citep{galen-mt-mqm-2026}, and
exactly what our own keyness punch list (\S\ref{sec:punch}) ranks. A single frequency axis cannot
surface these words; the medical-vs-non-medical contrast can, because it keys on a word being
\textit{at home in both} of Galen's registers rather than on raw scarcity. This is where the second
axis earns its keep --- not in re-finding the rare terms a rarity metric already flags, but in
isolating the population a rarity metric is blind to.

We are deliberate about this quadrant's status --- its name is \texttt{polysemy\_risk}, not a verdict.
That ordinary words systematically acquire specialized, register-bound senses in ancient technical
medicine is well established, if in the sister Latin tradition~\citep[ch.~3]{langslow-2000}; such
terms are, by their nature, the ones a reader can render in a plausible everyday sense while missing
the technical one. This is a claim about the \textit{structure of technical vocabulary}, not a finding
of the MT measurement we build on --- that measurement in fact reports the opposite for its own error
data, finding failure concentrated in \textit{rare} terminology while common technical terms translate
cleanly~\citep{galen-mt-mqm-2026}. Our claim is therefore not that these terms are hard for MT ---
they may not be --- but that a cheap, annotation-free frequency contrast \textbf{partitions
candidate-hard vocabulary into two mechanistically distinct classes}: rarity-driven difficulty, which
the existing predictor already captures, and sense-loading, which it cannot see. Which class actually
degrades machine translation, and by how much, is the open question \S\ref{sec:conclusion} proposes
to test.

\paragraph{The domain quadrant is the control.}
Its members --- σφυγμός, δραχμή, λίτρα, ῥίζα, ὄξος, κηρός,
σμύρνα, χαλβάνη, ὀπίον, ῥητίνη, ἀμμωνιακός --- are materia medica and units: rare against general
Greek, but ordinary within medicine. Whole-corpus keyness would flag these as ``hard''; the four-way
correctly files them as domain vocabulary a glossary resolves, a different and easier problem than
the polysemy set. The worked case is ἀμμωνιακός at medical 1.40 / non-medical 0.00.

The practical payoff is that the classification tells an MT researcher \textit{which kind} of hard a term
is, and the two kinds want different fixes --- a term glossary for the domain quadrant,
sense-disambiguation for the polysemy quadrant.

\section{The Punch List}
\label{sec:punch}

The human-facing deliverable is a punch list of \textbf{1{,}433} candidate terms passing the thin-gloss LSJ
filter, split by \texttt{lsj\_tier} into \textbf{absent} (734, no LSJ entry) and \textbf{thin} (699, $\le$1 sense line) ---
reported separately because they are different-quality signals: LSJ is comprehensive, so a genuine
rare medical term usually has a thin entry, whereas the absent bucket also collects lemmatizer
artifacts and accent mismatches. Each record carries the lemma, both frequencies, keyness, LSJ sense
count, the top example work URNs, and one attested example sentence, with the identification field
left empty --- the scaffold a Phase~3 fills.

As a cross-validation check, \textbf{97.1\%} of punch-list terms (1{,}391 / 1{,}433) fall in the
\texttt{genuinely\_hard} quadrant of the independent four-way classification, confirming the punch list is not
contaminated by domain vocabulary.

\section{Limitations}
\label{sec:limitations}

We state these plainly; under-claiming is the credibility of a resource paper.

\begin{itemize}[leftmargin=*,itemsep=3pt]
\item \textbf{Lemmatizer error is reported, not patched.} Three systematic classes remain and are handed to
  the reviewer rather than silently corrected: \textbf{fragmentation} (one headword split across variant
  lemmas, e.g.\ φλεβά/φλείβ $\rightarrow$ φλέψ, which distorts ranking by splitting a term's frequency);
  \textbf{consistently-wrong lemmas of irregular nouns} (νεύρος for νεῦρον), which the automated QC gates
  cannot catch by construction; and one such case surfaced only by manual audit (δεικέω $\rightarrow$ δείκνυμι).
  Neural-lemmatizer hallucinations --- non-words with no LSJ entry and no near-neighbor surface form
  anywhere in the corpus --- are auto-dropped (86, e.g.\ γαήτσρ); this is a ``not a Greek word''
  judgment, not a lexical one. The rare-stratum gold standard (\S\ref{sec:accuracy}) quantifies
  these classes at scale: of 697 mismatches on LSJ-cited loci, 411 are identity passthrough, 135
  garbled non-words, 136 wrong-lemma substitutions, and 15 variant-headword conventions.
\item \textbf{The machine-adjudicated accuracy estimate is not a human check.} Two language
  models passed the salted-control gate and agreed closely, but they share training lineage and
  are not Greek philologists; their estimate (\S\ref{sec:accuracy}) is reported under its own
  label and the human two-reader protocol stays planned.
\item \textbf{Census caveats are inherited:} the volume approximation is unverified, and category tags are
  author-assigned and contestable.
\item \textbf{One partial source} is documented rather than hidden: tlg035 (grc) is a 406-token Oribasius
  excerpt, not an extraction failure.
\item \textbf{Thin LSJ gloss $\neq$ medically undefined.} The thin-gloss filter is an approximate heuristic for
  ``under-described,'' not a claim that a term lacks a definition; we state it as such.
\end{itemize}

\section{Availability \& Reproducibility}
\label{sec:availability}

Census:~\citep{pepperman-2026-galen-census}, DOI
\url{https://doi.org/10.5281/zenodo.21561962}. Code: MIT. Data: CC-BY-4.0.
Underlying Greek texts:~\citep{first1kgreek}. The lemmatizer model and all library versions are pinned;
every stage is config-driven and re-runnable over a different subcorpus list, so the classification
can be reproduced, audited, or retargeted.

\section{Conclusion \& Future Work}
\label{sec:conclusion}

We have built a measured map of where Galen's translation-blocking vocabulary concentrates, and
shown that a cheap frequency signal partitions that vocabulary by the \textit{kind} of difficulty it
poses, isolating a common, sense-loaded class that rarity-based predictors miss by construction.

\paragraph{A 2026 stability probe: a negative result, reported.}
Because the map's diagnostic claim is about translation, we pre-registered and ran a cheap test of
it against 2026 models (\texttt{mt2026\_consistency\_probe\_spec.md}, with the data). Two frontier
models (\texttt{claude-fable-5-1}, \texttt{claude-opus-5-5}) translated the same 40 passages
blind --- 20 each from \textit{De temperamentis} and \textit{De compositione medicamentorum per
genera}, stratified by punch-list density --- and reported the rendering they used for 153
punch-list tokens and 153 matched control tokens (NOUN/ADJ in the same passages, not on the list,
LSJ sense count $\ge 3$). We measure \emph{stability}, not correctness: how often the two renderings
differ after an identical normalization. They differed on \textbf{10.5\%} of punch-list targets
and \textbf{9.8\%} of controls (difference 0.65 points, 95\% CI $-6.1$ to $+7.4$, $p = 0.85$);
restricted to the \texttt{thin} tier, 8.6\% against 9.8\%. Passage-level disagreement did not
track punch-list density (Spearman $\rho = 0.12$, 95\% CI $-0.24$ to $0.45$, $n = 36$). Two
qualifications: the two models share a developer and their agreement is therefore optimistic, and
the \texttt{absent} tier of the punch list is visibly contaminated by lemmatizer artifacts of
ordinary words (κρᾶσι for κρᾶσις, μυύς for μῦς), which the readers of \S\ref{sec:accuracy} would
have caught. A follow-up pass measured the contamination: an independent lexicon-first lemmatizer
(Open Greek's Dilemma, November 2026) re-lemmatizes 625 of the 734 \texttt{absent}-tier lemmas to
LSJ headwords, 489 of them already key lemmas of the map; two blind model readers, each catching all
59 salted controls, call \textbf{83.4\%} and \textbf{83.8\%} of the tier artifacts of another
headword ($\kappa = 0.88$); the 108 both call real headwords are Galen's pharmacological authorities
(38 names) and medical vocabulary with no entry in the Perseus LSJ (70: κακοχυμία, λειποθυμία,
περιπνευμονικός, καρωτίς, κνίκος \ldots). The \texttt{thin} tier is 88\% headword-stable under the
second lemmatizer. The v3.0 release carries these as attributed columns and re-labels the tiers; the
list itself is not patched. What the probe does show is where 2026 renderings \emph{do} wobble: on sense-loaded
ordinary words (πλῆθος, καιρός, τροφή, πολύς) and on units and preparations with competing
scholarly conventions (\textit{cotylae}/\textit{kotylai}, \textit{sextarii}/\textit{xestai},
\textit{dyscrasia}/\textit{bad-mixture}), not on the rare technical lexicon, which both models
rendered identically (litharge, cerate, colophony, omphacitis). On this evidence the punch list
remains a map of what the lexicon under-describes; it is not, for current models, a map of where
rendering is unstable. A second pre-registered probe tested the obvious alternative, the polysemy
quadrant of \S\ref{sec:results} (159 targets) against the medical-only quadrant (151), with three model
families and two independent runs each: polysemous words were \emph{not} rendered less stably
(within-model disagreement 8.8--11.9\% against 7.9--14.6\%; every difference pointed the other way, none
significant). What polysemy produces is context-dependent rendering of the same lemma across passages,
which is correct behaviour, not noise; what the medical-only lexicon produces is translator
\emph{uncertainty} (self-flagged in 2.7--11.3\% of its targets against 0--6.3\%, in all six runs) and,
on the ten punch-list terms we constrained in a third run, the highest within-model instability in their
passages (15--27\% against 6--11\%). That third run also showed that a glossary built mechanically from
LSJ's first glosses is obeyed at 100\% and judged worse by blind readers in three of four decided pairs:
the terminology a translator needs is a list of equivalents chosen by someone who knows the
pharmacopoeia, not a dictionary's definitions, and \citet{galen-mt-mqm-2026}'s rarity effect is, on this
evidence, an effect of the technical lexicon after all.

\paragraph{Phase 3}
compiles, for the punch-list terms, \textit{attributed} identifications already published in the
scholarly record (e.g.\ André on plant names, Beck's Dioscorides concordance, Berendes, LSJ) ---
gathering and citing, never adjudicating. A v0 of this list --- the full attested concordance,
K\"uhn loci, the LSJ entry text, the term's footprint across Hippocrates, Dioscorides, Oribasius,
Aetius, Paul and Aretaeus in the same-model reference corpus, and the machine QC flags --- ships
with this version; glosses and the printed reference works are the v1 that needs a human.

\paragraph{The demonstration that would close the loop:}
define ${\sim}$10 polysemy-quadrant terms, supply them as
constrained vocabulary to an MT system, re-translate a passage, and measure the MQM quality delta
against the $r = -.97$ predictor. That turns the punch list from a spreadsheet into evidence that
naming the vocabulary moves the number.

\paragraph{A second use: adjudicating word-sense claims with corpus evidence.}
The same structure that separates rarity-driven from sense-loaded difficulty also answers a question
that recurs outside MT --- whether a given Galenic term actually carries the special sense a reader
wants to assign it. Because every count is URN-keyed and partitioned by subcorpus, the map reports,
for any lemma, how often it occurs and where: ordinary across Galen's registers, or concentrated and
technical. That is precisely the evidence a philologist marshals by hand against the over-reading of
individual words --- a common term pressed to mean something esoteric --- supplied here at corpus
scale with a fixed, auditable provenance. We note this as a \textit{use}, not a result: the instrument
stays neutral about what any term means, but it makes the frequency-and-distribution facts that
constrain such claims cheap to retrieve.

\section*{Acknowledgments}

The author thanks M\'arton Kardos for generous correspondence about odyCy early in this
project, which helped set the direction of the accuracy measurements in \S\ref{sec:accuracy}.
Thanks also to the Open Greek and Latin project, whose First1KGreek digitizations make
corpus-scale work on Galen possible.

\bibliography{references}

\end{document}
